\section{Introducción y Objetivos}

\subsection{Introducción}
El sector comercial de productos y alimentos balanceados para mascotas en el Ecuador ha experimentado un crecimiento acelerado en los últimos años, impulsado por una mayor concientización sobre el bienestar animal y la tenencia responsable. Sin embargo, establecimientos comerciales de tamaño mediano como ``PetShop Huellitas \& Snacks'' se enfrentan a desafíos significativos al gestionar sus existencias y canales de atención de manera tradicional y manual. La falta de automatización y de sincronía entre el almacén físico y las solicitudes de pedidos genera ineficiencias operativas y constantes quiebres de stock. 

Ante este escenario, la implementación de un sistema de software enfocado a la automatización de procesos operativos se presenta como una solución fundamental para optimizar la cadena de valor del negocio y garantizar su supervivencia frente al ingreso de competidores de gran escala. Este informe presenta el análisis y diseño detallado de la solución informática denominada ``PetFood-Manager'', estructurando sus requerimientos, diagramas de interacción y diagramas estructurales para fundamentar un desarrollo de software robusto, transaccional y escalable.

\subsection{Objetivos del Proyecto}

\subsubsection{Objetivo General}
Diseñar el modelo de sistema informático ``PetFood-Manager'' para la tienda de venta de comida de mascotas ``PetShop Huellitas \& Snacks'', con el fin de automatizar el control de inventario en tiempo real, agilizar el procesamiento de pedidos de clientes y sistematizar el reabastecimiento con los proveedores bajo una estrategia de supervivencia eficiente.

\subsubsection{Objetivos Específicos}
\begin{itemize}
    \item Contextualizar el flujo de información de la tienda mediante el modelado del diagrama de contexto y el diagrama de flujo de datos (DFD) de nivel 0, delimitando claramente las interacciones con clientes, proveedores y entes externos.
    \item Modelar el comportamiento dinámico y la estructura de clases del sistema utilizando la notación UML, detallando la interacción lógica de los componentes de software mediante diagramas de secuencia, actividades y clases.
    \item Definir la arquitectura física y de componentes del sistema considerando un despliegue distribuido en la nube, garantizando requerimientos no funcionales de seguridad, integridad y alta disponibilidad de la información.
\end{itemize}